\documentclass[11pt, a4paper]{article}
\usepackage[utf8]{inputenc}
\usepackage{amsmath, amssymb}
\usepackage[round, authoryear]{natbib}
\usepackage{geometry}
\geometry{margin=1in}

\begin{document}

\section{Theoretical Architecture of the Life-Optimisation Markov Decision Process}

To operationalise human life decisions within a formal reinforcement learning and dynamic
programming framework, the environment must be structured as a Markov Decision Process (MDP).
The fundamental requirement of an MDP is the Markov property: the future state of the system
must depend solely upon the present state and the current action taken, rendering the precise
historical sequence of prior states mathematically irrelevant to forward-looking optimisation.
Human life is inherently path-dependent, which initially appears to violate this memoryless
assumption. However, by carefully engineering the state space to act as \emph{sufficient
statistics} of history---compressing years of cumulative actions into current structural
constraints---the Markov property can be restored. Furthermore, traditional exponential
discounting fails to capture well-documented psychological realities of human choice. This
section justifies a six-dimensional state representation and a modified quasi-hyperbolic
Bellman equation featuring a backward-looking habit stock.

%─────────────────────────────────────────────────────────────────────────────
\subsection{Validating the Six-Dimensional Markov State Space}
%─────────────────────────────────────────────────────────────────────────────

The proposed life-optimisation MDP utilises a bounded six-dimensional state space vector,
$S_t$. Each dimension represents an accumulated integral or temporal derivative of past
behaviour, successfully satisfying the Markov property by summarising all necessary historical
context required to calculate transition probabilities and expected future utility.

\subsubsection{Financial Pressure (\textit{critical} $\rightarrow$ \textit{comfortable})}

Financial pressure encapsulates the aggregate monetary history of the agent, functioning as a
continuous state variable ranging from a critical liquidity crunch to comfortable abundance.
Within a dynamic programming framework, current wealth serves as the definitive state variable
that dictates feasible action spaces, particularly under liquidity constraints
\citep{Deaton1991}. If an agent's current financial pressure state is known, the specific
historical sequence of consumption and income generation provides no additional predictive
power regarding their future purchasing capabilities. Furthermore, behavioural economics
demonstrates that financial pressure transcends mere budget constraints; scarcity acts as a
cognitive bandwidth tax that fundamentally alters an individual's decision-making architecture
\citep{Mullainathan2013}. Thus, financial pressure operates as a crucial Markov state that
parametrises both the physical action space and the internal cognitive processing capacity of
the agent.

\subsubsection{Energy Level (\textit{exhausted} $\rightarrow$ \textit{peak})}

The energy level dimension, ranging from \textit{exhausted} to \textit{peak}, compresses the
agent's immediate physiological and metabolic history (including sleep, exertion, and diet)
into a single, exhaustible resource budget. To predict the transition probability of an agent
successfully executing a high-effort action rather than defaulting to a low-effort alternative,
the optimiser requires only the current depletion level. This is deeply grounded in the
psychological model of ego depletion, which posits that self-control and volition draw upon a
limited, quantifiable physiological resource \citep{Baumeister1998}. When this state variable
approaches \textit{exhausted}, the probability of successfully transitioning to states
requiring high executive function approaches zero. Biologically, this cumulative wear-and-tear
is formalised as allostatic load \citep{McEwen1998}, providing a robust physiological
justification for treating energy as a strict Markovian state variable governing action
efficacy.

\subsubsection{Emotional State (\textit{stressed} $\rightarrow$ \textit{excited})}

Emotional state operates as a transient parameter that temporarily rewrites the agent's reward
function and risk tolerance, scaling from \textit{stressed} and anxious to \textit{excited}
and calm. \citet{Loewenstein1996} illustrates that such ``visceral factors'' act as profound
state variables capable of overriding long-term cognitive goals; extreme emotional states
fundamentally alter the expected subjective value of future rewards. Moreover, prospect theory
dictates that human decisions under risk are evaluated relative to a current reference point
rather than absolute final outcomes \citep{Kahneman1979}. The emotional state serves as this
shifting reference point. Knowing the sequence of past traumas is less mathematically relevant
to predicting immediate behavioural responses than knowing the precise, current visceral state
of the agent at time $t$.

\subsubsection{Social Connectedness (\textit{isolated} $\rightarrow$ \textit{thriving})}

Relationships require continuous maintenance, meaning social connectedness functions
identically to a depreciating asset. Ranging from \textit{isolated} to \textit{thriving},
this dimension represents the current accumulated stock of social capital. Sociological
research extensively documents how accumulated social capital serves as a predictive reservoir
of future utility, emotional buffering, and opportunity \citep{Putnam2000}. By explicitly
modelling social connections within a Bellman framework, \citet{Glaeser2002} demonstrate that
social capital is a valid state variable that demands time investment; actions increase the
stock, while inaction leads to steady algorithmic decay. Consequently, the current state of
social connectedness is entirely sufficient to calculate the probability of receiving social
rewards in the next time step.

\subsubsection{Skill Momentum (\textit{declining} $\rightarrow$ \textit{accelerating})}

Traditional models often utilise static skill level as a state variable; however, including
\emph{skill momentum} (the temporal derivative of human capital, ranging from
\textit{declining} to \textit{accelerating}) provides a significantly more robust Markovian
representation in environments with high friction. Foundational economic theory establishes
human capital as an accumulated investment stock that alters the transition probabilities of
future earnings \citep{Becker1962}. By augmenting this with momentum, the MDP captures the
psychological reality of learning. Behavioural momentum theory quantitatively proves that an
organism's resistance to abandoning a task in the face of friction is directly proportional to
their current momentum state \citep{Nevin1992}. Therefore, if momentum is
\textit{accelerating}, the transition probability of project abandonment remains low
regardless of distant past failures, preserving the memoryless requirement of the MDP.

\subsubsection{Life Stage (\textit{student} $\rightarrow$ \textit{senior})}

Human life is fundamentally a finite-horizon problem. The expected value of long-term
investments inherently changes depending on the time remaining to amortise those investments.
In reinforcement learning, introducing a time-step or phase variable is the standard
algorithmic mechanism used to convert a non-stationary, finite-horizon environment into a
stationary, Markovian one \citep{Sutton2018}. This mirrors the life-cycle hypothesis in
economics, which utilises life stage (from \textit{student} to \textit{senior}) as the master
state variable dictating macro-shifts in rational policy, moving an agent from heavy borrowing
and human capital acquisition in early stages toward accumulation and eventual decumulation in
later stages \citep{Modigliani1954}.

%─────────────────────────────────────────────────────────────────────────────
\subsection{Quasi-Hyperbolic Discounting with Backward-Looking Habit Stocks}
%─────────────────────────────────────────────────────────────────────────────

Standard Bellman equations rely on exponential discounting ($\delta^{t}$) to ensure time
consistency. However, empirical human behaviour routinely violates this assumption due to
myopia \citep{Strotz1955} and present bias. To construct an ecologically valid life-optimiser,
the MDP must discard strict exponential discounting in favour of a quasi-hyperbolic
($\beta$-$\delta$) discount function combined with a backward-looking habit stock. The
modified action-value function is formalised as:

\begin{equation}
  Q(s,a) \;=\; h(\text{past}) \;+\; R(s,a) \;+\; \beta \cdot \delta
  \cdot \sum_{s'} P(s'\mid s,a)\cdot V(s')
  \label{eq:quasi-hyperbolic-Q}
\end{equation}

\noindent
where $R(s,a)$ is the immediate reward and the summation represents the expected future value
weighted by transition probabilities $P(s'\mid s,a)$.

\paragraph{The long-run discount factor $\delta$.}
The standard exponential discount factor $\delta \in (0,1)$ captures rational, consistent
long-term time preference. Its magnitude is inferred from personality traits: higher
conscientiousness and rationality yield higher $\delta$, consistent with studies on
delay-of-gratification and executive function \citep{Sutton2018}.

\paragraph{The present-bias coefficient $\beta$.}
The critical addition is the parameter $\beta \in (0,1]$, which uniformly down-weights all
future periods relative to the present moment, mathematically operationalising quasi-hyperbolic
discounting. The original $\beta$-$\delta$ functional form was introduced by
\citet{Phelps1968} and later applied to individual decision-making in a Bellman framework by
\citet{Laibson1997}. When $\beta < 1$, the model captures the present bias consistently
observed in behavioural economics \citep{Frederick2002}. This framework allows the optimiser
to model an agent who plays an intra-personal game against their future selves, and to
determine whether the agent is computationally \textit{na\"ive} (unaware of their own future
bias) or \textit{sophisticated} (aware of their bias and capable of selecting commitment
actions accordingly) \citep{ODonoghue1999}.

\paragraph{The habit stock $h(\text{past})$.}
The term $h(\text{past})$ introduces a backward-looking utility modifier, representing the
exponentially decayed memory of past rewards:

\begin{equation}
  h(\text{past}) \;=\; \sum_{i=1}^{N} \alpha^{\,\Delta t_i} \cdot R_{t-i}
  \label{eq:habit-stock}
\end{equation}

\noindent
where $\alpha \in (0,1)$ is the memory decay rate and $\Delta t_i$ is the age of observation
$i$ (in hours). This formulation aligns with the mathematics of habit formation established by
\citet{Ryder1973}, where current utility is a function of the deviation from past consumption
levels. It further accommodates regret theory, under which the memory of foregone
counterfactual states explicitly alters the valence of the current reward \citep{Loomes1982},
and case-based decision theory, where historical payoffs are dynamically weighted by similarity
to the current situation \citep{Gilboa1995}. By compressing historical vectors into the scalar
$h(\text{past})$, the architecture incorporates memory and regret without violating the
Markov property, resulting in a mathematically rigorous yet psychologically realistic
optimisation engine.

%─────────────────────────────────────────────────────────────────────────────
\newpage
\begin{thebibliography}{99}

\bibitem[Baumeister et al.(1998)]{Baumeister1998}
  Baumeister, R.F., Bratslavsky, E., Muraven, M.\ \& Tice, D.M. (1998).
  Ego Depletion: Is the Active Self a Limited Resource?
  \textit{Journal of Personality and Social Psychology}, \textbf{74}(5), 1252--1265.

\bibitem[Becker(1962)]{Becker1962}
  Becker, G.S. (1962).
  Investment in Human Capital: A Theoretical Analysis.
  \textit{Journal of Political Economy}, \textbf{70}(5, Part~2), 9--49.

\bibitem[Deaton(1991)]{Deaton1991}
  Deaton, A. (1991).
  Saving and Liquidity Constraints.
  \textit{Econometrica}, \textbf{59}(5), 1221--1248.

\bibitem[Frederick, Loewenstein and O'Donoghue(2002)]{Frederick2002}
  Frederick, S., Loewenstein, G.\ \& O'Donoghue, T. (2002).
  Time Discounting and Time Preference: A Critical Review.
  \textit{Journal of Economic Literature}, \textbf{40}(2), 351--401.

\bibitem[Gilboa and Schmeidler(1995)]{Gilboa1995}
  Gilboa, I.\ \& Schmeidler, D. (1995).
  Case-Based Decision Theory.
  \textit{The Quarterly Journal of Economics}, \textbf{110}(3), 605--639.

\bibitem[Glaeser, Laibson and Sacerdote(2002)]{Glaeser2002}
  Glaeser, E.L., Laibson, D.\ \& Sacerdote, B. (2002).
  An Economic Approach to Social Capital.
  \textit{The Economic Journal}, \textbf{112}(483), F437--F458.

\bibitem[Kahneman and Tversky(1979)]{Kahneman1979}
  Kahneman, D.\ \& Tversky, A. (1979).
  Prospect Theory: An Analysis of Decision under Risk.
  \textit{Econometrica}, \textbf{47}(2), 263--291.

\bibitem[Laibson(1997)]{Laibson1997}
  Laibson, D. (1997).
  Golden Eggs and Hyperbolic Discounting.
  \textit{The Quarterly Journal of Economics}, \textbf{112}(2), 443--478.

\bibitem[Loewenstein(1996)]{Loewenstein1996}
  Loewenstein, G. (1996).
  Out of Control: Visceral Influences on Behavior.
  \textit{Organizational Behavior and Human Decision Processes},
  \textbf{65}(3), 272--292.

\bibitem[Loomes and Sugden(1982)]{Loomes1982}
  Loomes, G.\ \& Sugden, R. (1982).
  Regret Theory: An Alternative Theory of Rational Choice Under Uncertainty.
  \textit{The Economic Journal}, \textbf{92}(368), 805--824.

\bibitem[McEwen(1998)]{McEwen1998}
  McEwen, B.S. (1998).
  Stress, Adaptation, and Disease: Allostatic Load and Allostatic State.
  \textit{Annals of the New York Academy of Sciences}, \textbf{840}, 33--44.

\bibitem[Modigliani and Brumberg(1954)]{Modigliani1954}
  Modigliani, F.\ \& Brumberg, R. (1954).
  Utility Analysis and the Consumption Function: An Interpretation of
  Cross-Section Data. In K.~Kurihara (Ed.),
  \textit{Post-Keynesian Economics} (pp.~388--436).
  New Brunswick, NJ: Rutgers University Press.

\bibitem[Mullainathan and Shafir(2013)]{Mullainathan2013}
  Mullainathan, S.\ \& Shafir, E. (2013).
  \textit{Scarcity: Why Having Too Little Means So Much}.
  New York: Henry Holt and Company.

\bibitem[Nevin(1992)]{Nevin1992}
  Nevin, J.A. (1992).
  An Integrative Model for the Study of Behavioral Momentum.
  \textit{Journal of the Experimental Analysis of Behavior},
  \textbf{57}(3), 301--316.

\bibitem[O'Donoghue and Rabin(1999)]{ODonoghue1999}
  O'Donoghue, T.\ \& Rabin, M. (1999).
  Doing It Now or Later.
  \textit{American Economic Review}, \textbf{89}(1), 103--124.

\bibitem[Phelps and Pollak(1968)]{Phelps1968}
  Phelps, E.S.\ \& Pollak, R.A. (1968).
  On Second-Best National Saving and Game-Equilibrium Growth.
  \textit{The Review of Economic Studies}, \textbf{35}(2), 185--199.

\bibitem[Putnam(2000)]{Putnam2000}
  Putnam, R.D. (2000).
  \textit{Bowling Alone: The Collapse and Revival of American Community}.
  New York: Simon \& Schuster.

\bibitem[Ryder and Heal(1973)]{Ryder1973}
  Ryder, H.E.\ \& Heal, G.M. (1973).
  Optimal Growth with Intertemporally Dependent Preferences.
  \textit{The Review of Economic Studies}, \textbf{40}(1), 1--31.

\bibitem[Strotz(1955)]{Strotz1955}
  Strotz, R.H. (1955).
  Myopia and Inconsistency in Dynamic Utility Maximization.
  \textit{The Review of Economic Studies}, \textbf{23}(3), 165--180.

\bibitem[Sutton and Barto(2018)]{Sutton2018}
  Sutton, R.S.\ \& Barto, A.G. (2018).
  \textit{Reinforcement Learning: An Introduction} (2nd~ed.).
  Cambridge, MA: MIT Press.

\end{thebibliography}

\end{document}
